% Audit de veracite des citations de rapport de jury - 2026-09-20
% Bloc 4 : citation jury 2023 (numeration / hexadecimal / BCD) introuvable dans
% le corpus ; remplacee par le verbatim exact du Rapport_Jury_2023 (l.463).
\expandafter\ifx\csname CAPETINCLUDEMODE\endcsname\relax
\documentclass[11pt,a4paper]{article}
% preamble.tex — packages et commandes communes a toutes les fiches CAPET SII IE
\input{glyphtounicode}
\pdfgentounicode=1
\usepackage{cmap}
\usepackage[T1]{fontenc}
\usepackage[utf8]{inputenc}
\usepackage[french]{babel}
% babel[french] active a \begin{document} des raccourcis de ponctuation (:;!?)
% qui cassent la syntaxe to[...] de circuitikz (erreur "+ or - expected") ;
% on les desactive, ce qui ne change que l'espacement fin avant ces signes.
\AtBeginDocument{\shorthandoff{:;!?}}
\usepackage[a4paper, top=1.6cm, bottom=1.6cm, left=1.4cm, right=1.4cm]{geometry}
\setlength{\headheight}{14.5pt}
\usepackage{amsmath, amssymb, amsfonts}
\usepackage{siunitx}
\sisetup{locale=FR, output-decimal-marker={,}, inter-unit-product=\ensuremath{{}\cdot{}},
         group-minimum-digits=4}
% Unites non fournies par siunitx, declarees ici une fois pour toutes (fiches + exercices).
% Les gardes \ifdefined evitent tout conflit avec une redeclaration locale restee dans un
% fichier : ne pas retirer les gardes, ne pas redeclarer ces unites dans les fichiers.
\DeclareSIUnit{\jour}{j}
\ifdefined\var\else\DeclareSIUnit{\var}{var}\fi   % puissance reactive
\ifdefined\vad\else\DeclareSIUnit{\vad}{VAD}\fi   % puissance deformante
\ifdefined\tr\else\DeclareSIUnit{\tr}{tr}\fi      % tour (tr/min)
\usepackage{array, booktabs, multirow, tabularx}
\usepackage{xcolor}
\usepackage[most, breakable]{tcolorbox}
\usepackage{multicol}
\usepackage{enumitem}
\usepackage{mathtools}
\usepackage{lmodern}
% Refonte visuelle 2026 (direction "moderne coloré") : Lato pour le texte,
% maths sans-serif assorties via newtxsf (charge apres amsmath/amssymb).
\usepackage[default]{lato}
\usepackage{newtxsf}
\usepackage{microtype}
\usepackage{fancyhdr}
\usepackage{tikz}
\usetikzlibrary{positioning,arrows.meta}
\usepackage{circuitikz}
\usepackage{lastpage}

% ── Couleurs par domaine (palette modernisée 2026) ────────────────────
\definecolor{colElec}{RGB}{29, 78, 216}        % Électrotechnique — bleu vif
\definecolor{colElecLight}{RGB}{219, 234, 254} % fond bleu clair
\definecolor{colAuto}{RGB}{21, 128, 61}        % Automatique — vert
\definecolor{colMeca}{RGB}{194, 65, 12}        % Mécanique — orange
\definecolor{colTherm}{RGB}{185, 28, 28}       % Thermique — rouge
\definecolor{colQT}{RGB}{232, 244, 233}        % (hérité, plus utilisé par questiontype)
\definecolor{colJury}{RGB}{255, 248, 225}      % fond jury (crème)
\definecolor{colJuryFrame}{RGB}{202, 138, 4}   % bordure/bandeau jury (ambre)
\definecolor{colGrayLight}{RGB}{245, 245, 247} % fond gris neutre
\definecolor{colSeparator}{RGB}{180, 180, 180} % gris séparateur
\definecolor{colExo}{RGB}{8, 145, 178}          % Banque d'exercices — cyan (distinct des fiches)
% ── Encre et gris de la refonte ───────────────────────────────────────
\definecolor{colInk}{RGB}{23, 32, 42}          % texte courant (quasi-noir)
\definecolor{colSlate}{RGB}{83, 98, 118}       % gris ardoise (méta, en-têtes)
\definecolor{colRule}{RGB}{203, 213, 225}      % filets clairs
\definecolor{colJuryBg}{RGB}{255, 248, 225}    % fond crème jury
\definecolor{colPresence}{RGB}{83, 98, 118}    % bloc 5 (présence dans les sujets)
\AtBeginDocument{\color{colInk}}

% ── Style encadré exercice + corrigé (banque d'exercices) ─────────────
\tcbset{
  exoBox/.style={
    breakable,
    enhanced,
    colback=colExo!8!white,
    colframe=colExo,
    colbacktitle=colExo,
    coltitle=white,
    fonttitle=\bfseries\small,
    left=4pt, right=4pt, top=3pt, bottom=3pt,
    before skip=5pt, after skip=5pt,
  }
}

% ── Bascule enonce / corrige (banque d'exercices) ─────────────────────
% \CAPETCORRIGE est defini en ligne de commande (pdflatex "\def\CAPETCORRIGE{1}\input{...}")
% pour compiler la variante corriges d'un meme fichier source ; par defaut (non defini),
% la compilation produit la variante enonces. \CAPETSHOW{enonce}{corrige} choisit la branche.
%
% TROIS MODES (refonte 2026-09-20, manuel de revision) :
%   - aucun drapeau   -> imprime l'enonce            (variante enonces, INCHANGE)
%   - \CAPETCORRIGE   -> imprime le corrige          (variante corriges, INCHANGE)
%   - \CAPETMANUEL    -> imprime l'enonce ET met le corrige de cote pour l'annexe
%
% Le mode manuel stocke chaque corrige dans une macro (\csgdef), sans ecriture
% fichier ni re-tokenisation : le contenu (circuitikz, accents, %) traverse intact.
% L'annexe les rejoue ensuite dans l'ordre via \CAPETANNEXE.
\usepackage{etoolbox}

\newcounter{CAPETexo}          % numero continu d'exercice dans le manuel
\setcounter{CAPETexo}{0}
\newcounter{CAPETannexe}       % compteur de relecture en annexe

% Titre courant de la notion en cours (pose par \CAPETNOTION), sert d'intertitre annexe
\newcommand{\CAPETnotioncode}{}
\newcommand{\CAPETnotiontitre}{}

% \CAPETNOTION{code}{titre} : ouvre une notion dans le manuel (memorise le libelle,
% et marque le rang du premier exercice de la notion pour l'intertitre d'annexe).
\newcommand{\CAPETNOTION}[2]{%
  \renewcommand{\CAPETnotioncode}{#1}%
  \renewcommand{\CAPETnotiontitre}{#2}%
  \csgdef{CAPETnotionstart\the\numexpr\value{CAPETexo}+1\relax}{#1 -- #2}%
}

% ── \degree robuste ───────────────────────────────────────────────────
% Certains fichiers d'exercices utilisent \degree. Ni le preambule ni siunitx
% ne le definissent, et \providecommand local ne suffit pas : sous hyperref,
% \degree existe deja comme commande factice (\ERROR), donc \providecommand
% s'abstient et l'appel echoue a la compilation du manuel (erreur invisible
% en compilation isolee). On le (re)definit donc inconditionnellement ici,
% apres le chargement des packages, pour tout le corpus.
\AtBeginDocument{\providecommand{\degree}{}\renewcommand{\degree}{\ensuremath{^\circ}}}

% Libelle decoratif : les fichiers d'exercices utilisent aussi \CAPETSHOW dans
% leur \fancyhead et leur bloc de titre, pour afficher "(enonces)" ou "(corriges)".
% Ces appels-la ne sont PAS des exercices : en mode manuel ils ne doivent ni
% consommer un numero, ni alimenter l'annexe. On les reconnait a leur premier
% argument, strictement egal au libelle ci-dessous.
\newcommand{\CAPETLIBELLEENONCES}{(\'enonc\'es)}

\ifdefined\CAPETMANUEL
  % --- Mode manuel : enonce en place, corrige differe vers l'annexe ---
  \newcommand{\CAPETSHOW}[2]{%
    \def\CAPETargun{#1}%
    \ifx\CAPETargun\CAPETLIBELLEENONCES
      #1%            appel decoratif : rendu tel quel, hors numerotation
    \else
      \CAPETSHOWEXO{#1}{#2}%
    \fi
  }
  \newcommand{\CAPETSHOWEXO}[2]{%
    \stepcounter{CAPETexo}%
    \csgdef{CAPETcorrige\theCAPETexo}{#2}%
    \edef\CAPETtmpnum{\theCAPETexo}%
    #1%
    \par\vspace{2pt}%
    {\footnotesize\itshape\color{colExo}%
     \hyperref[capet:cor:\CAPETtmpnum]{Corrige de l'exercice \CAPETtmpnum\ $\rightarrow$ p.~\pageref{capet:cor:\CAPETtmpnum}}%
     \label{capet:eno:\CAPETtmpnum}}%
  }
  % \CAPETANNEXE : rejoue tous les corriges captures, dans l'ordre, avec intertitres
  \newcommand{\CAPETANNEXE}{%
    \setcounter{CAPETannexe}{0}%
    \loop\ifnum\value{CAPETannexe}<\value{CAPETexo}%
      \stepcounter{CAPETannexe}%
      \edef\CAPETcurnum{\theCAPETannexe}%
      \ifcsname CAPETnotionstart\CAPETcurnum\endcsname
        \par\medskip
        {\normalsize\bfseries\color{colExo}\csuse{CAPETnotionstart\CAPETcurnum}}%
        \par\nobreak\smallskip
      \fi
      \begin{tcolorbox}[exoBox, breakable,
        title={Corrige de l'exercice \CAPETcurnum\
               \hfill\normalfont\footnotesize
               \hyperref[capet:eno:\CAPETcurnum]{enonce p.~\pageref{capet:eno:\CAPETcurnum}}}]%
      \label{capet:cor:\CAPETcurnum}%
      \csuse{CAPETcorrige\CAPETcurnum}%
      \end{tcolorbox}%
    \repeat
  }
\else
  % --- Modes historiques, comportement strictement inchange ---
  \newcommand{\CAPETSHOW}[2]{\ifdefined\CAPETCORRIGE #2\else #1\fi}
  \newcommand{\CAPETANNEXE}{}
\fi

% ── Difficulté 1 à 3 (étoiles pleines, amssymb) ───────────────────────
\newcommand{\difficulte}[1]{%
  \ifcase#1\relax
  \or $\bigstar$\;niveau 1/3%
  \or $\bigstar\bigstar$\;niveau 2/3%
  \or $\bigstar\bigstar\bigstar$\;niveau 3/3%
  \fi
}

% ── Style encadré extrait de document technique (banque d'exercices) ──
\tcbset{
  docTech/.style={
    breakable,
    enhanced,
    colback=colGrayLight,
    colframe=colSeparator,
    fonttitle=\bfseries\small\itshape,
    coltitle=black,
    colbacktitle=colGrayLight,
    left=4pt, right=4pt, top=3pt, bottom=3pt,
    before skip=4pt, after skip=4pt,
  }
}

% ── Style tcolorbox de section principale ────────────────────────────
% Refonte 2026 : bandeau titre coloré arrondi + fond très pâle, sans cadre.
\tcbset{
  ficheSection/.style={
    breakable, enhanced jigsaw,
    frame hidden, boxrule=0pt, arc=1.8mm,
    colback=#1!4!white,
    colbacktitle=#1,
    coltitle=white,
    fonttitle=\sffamily\bfseries,
    left=7pt, right=7pt, top=5pt, bottom=5pt,
    toptitle=3pt, bottomtitle=3pt,
    before skip=10pt, after skip=8pt,
  }
}

% ── Style encadré question type ──────────────────────────────────────
% Carte blanche, bandeau sombre, liseré très clair.
\tcbset{
  questiontype/.style={
    breakable, enhanced jigsaw,
    colback=white,
    frame hidden, boxrule=0pt, arc=1.8mm,
    colbacktitle=colElec!78!black,
    coltitle=white,
    fonttitle=\sffamily\bfseries\small,
    left=7pt, right=7pt, top=4pt, bottom=4pt,
    toptitle=2.5pt, bottomtitle=2.5pt,
    before skip=8pt, after skip=8pt,
    borderline={0.6pt}{0pt}{colElec!12!white},
  }
}

% ── Style encadré jury ───────────────────────────────────────────────
% Bandeau ambre, fond crème, sans cadre.
\tcbset{
  jury/.style={
    breakable, enhanced jigsaw,
    colback=colJuryBg,
    frame hidden, boxrule=0pt, arc=1.8mm,
    colbacktitle=colJuryFrame,
    coltitle=white,
    fonttitle=\sffamily\bfseries,
    left=7pt, right=7pt, top=4pt, bottom=4pt,
    toptitle=3pt, bottomtitle=3pt,
    before skip=10pt, after skip=8pt,
  }
}

% ── En-tête / pied de page communs (refonte 2026) ─────────────────────
% #1 code fiche, #2 titre court, #3 couleur domaine
\newcommand{\ficheHeaders}[3]{%
  \fancyhead[L]{\small\sffamily{\bfseries\textcolor{#3}{#1}}\hspace{0.7em}\textcolor{colSlate}{#2}}%
  \fancyhead[R]{\small\sffamily\textcolor{colSlate}{CAPET SII IE\, --- \,Session 2026}}%
  \fancyfoot[C]{\small\sffamily\textcolor{colSlate}{\thepage\,/\,\pageref{LastPage}}}%
  \fancyhead[C]{}\fancyfoot[L]{}\fancyfoot[R]{}%
}

% ── Bloc en-tête de fiche (BLOC 0, refonte 2026) ──────────────────────
% #1 code  #2 titre  #3 couleur domaine  #4 domaine  #5 ligne présence  #6 sous-thèmes
\newcommand{\ficheEntete}[6]{%
  \begin{tcolorbox}[enhanced, breakable,
    frame hidden, boxrule=0pt, arc=2.5mm,
    interior style={left color=#3, right color=#3!75!black},
    left=9pt, right=9pt, top=7pt, bottom=7pt,
    before skip=0pt, after skip=6pt]
    {\color{white}\sffamily\bfseries\LARGE #1\quad #2}\\[4pt]
    {\color{#3!22!white}\sffamily\small \MakeUppercase{#4}\quad\textbar\quad #5}\\[1pt]
    {\color{#3!35!white}\sffamily\footnotesize Sous-thèmes : #6}
  \end{tcolorbox}%
}

% ── Mise en page ─────────────────────────────────────────────────────
\pagestyle{fancy}
\fancyhf{}
\renewcommand{\headrulewidth}{0.5pt}
\renewcommand{\footrulewidth}{0pt}
\AtBeginDocument{%
  \renewcommand{\headrule}{{\color{colRule}\hrule width\headwidth height 0.5pt}}%
}

\setlist[itemize]{leftmargin=1.2em, itemsep=1pt, topsep=2pt, parsep=0pt}
\setlist[enumerate]{leftmargin=1.2em, itemsep=3pt, topsep=2pt, parsep=0pt}

\setlength{\parskip}{2pt}
\setlength{\columnsep}{10pt}

\def\CAPETBEGINFILE{\begin{document}}
\def\CAPETENDFILE{\end{document}}
\else
\def\CAPETBEGINFILE{}
\def\CAPETENDFILE{}
\fi

\ficheHeaders{T28}{Logique combinatoire}{colAuto}

\CAPETBEGINFILE

% -----------------------------------------------------------------------
%  BLOC EN-TETE
% -----------------------------------------------------------------------

\ficheEntete{T28}{Logique combinatoire}{colAuto}%
  {électronique numérique - Logique et codage}%
  {\textbf{Présence :} \textbf{3/13} sujets (implicite) - thème de \emph{culture}}%
  {algèbre de Boole - portes logiques - De Morgan - table de vérité - Karnaugh - circuits combinatoires - codage binaire / BCD / Gray}

% -----------------------------------------------------------------------
%  1. RAPPELS DE COURS
% -----------------------------------------------------------------------

\begin{tcolorbox}[
  ficheSection=colAuto,
  title={1. Rappels de cours},
]

% -- 1.1 Operateurs fondamentaux --------------------------------------
\textbf{1.1 Algèbre de Boole - opérateurs et théorèmes}

\smallskip
\begin{center}
\renewcommand{\arraystretch}{1.25}
\begin{tabular}{llll}
\toprule
\textbf{Opération} & \textbf{Notation} & \textbf{Identités notables} & \\
\midrule
ET (AND) & $A \cdot B$ & Idempotence : $A \cdot A = A$ & \\
OU (OR) & $A + B$ & Absorption : $A(A+B)=A$ & \\
NON (NOT) & $\overline{A}$ & Complémentarité : $A \cdot \overline{A} = 0$ & \\
\midrule
\multicolumn{4}{l}{\textbf{Théorèmes de De Morgan :}} \\
\multicolumn{4}{l}{$\overline{A \cdot B} = \overline{A} + \overline{B}$ \hspace{2em} $\overline{A + B} = \overline{A} \cdot \overline{B}$} \\
\bottomrule
\end{tabular}
\end{center}

\medskip\hrule\vspace{4pt}

% -- 1.2 Portes logiques -----------------------------------------------
\textbf{1.2 Portes logiques élémentaires (2 entrées)}

\begin{center}
\renewcommand{\arraystretch}{1.2}
\begin{tabular}{lll}
\toprule
\textbf{Porte} & \textbf{Expression} & \textbf{Remarque} \\
\midrule
AND & $S = A \cdot B$ & $S = 1$ ssi $A = B = 1$ \\
OR & $S = A + B$ & $S = 0$ ssi $A = B = 0$ \\
NOT & $S = \overline{A}$ & Inverseur \\
NAND & $S = \overline{A \cdot B}$ & Porte universelle \\
NOR & $S = \overline{A + B}$ & Porte universelle \\
XOR & $S = A \oplus B$ & $S = 1$ ssi $A \neq B$ \\
XNOR & $S = \overline{A \oplus B}$ & $S = 1$ ssi $A = B$ \\
\bottomrule
\end{tabular}
\end{center}

\textbf{Réalisation universelle :} toute fonction est réalisable avec des NAND seuls (ou NOR seuls).

\medskip\hrule\vspace{4pt}

% -- 1.3 Table de verite et formes canoniques --------------------------
\textbf{1.3 Table de vérité - formes canoniques}

Pour $n$ variables : $2^n$ lignes.
\begin{itemize}
  \item \textbf{SOP (Somme de mintermes)} : somme des mintermes pour lesquels $S = 1$
  \item \textbf{POS (Produit de maxtermes)} : produit des maxtermes pour lesquels $S = 0$
\end{itemize}

\textit{Exemple - demi-additionneur :}
\begin{center}
\renewcommand{\arraystretch}{1.1}
\begin{tabular}{cccc}
\toprule
$A$ & $B$ & $S = A \oplus B$ & $C = A \cdot B$ \\
\midrule
0 & 0 & 0 & 0 \\
0 & 1 & 1 & 0 \\
1 & 0 & 1 & 0 \\
1 & 1 & 0 & 1 \\
\bottomrule
\end{tabular}
\end{center}

\smallskip
\textbf{Schéma du demi-additionneur :}
\begin{center}
\begin{tikzpicture}[>=Stealth, font=\small, scale=0.95]
  \node (A) at (0,1.6) {$A$};
  \node (B) at (0,0.4) {$B$};
  \node[draw, minimum width=13mm, minimum height=9mm] (xor) at (3.2,1.5) {XOR};
  \node[draw, minimum width=13mm, minimum height=9mm] (and) at (3.2,0.5) {AND};
  \draw[->] (A) -| (2.4,1.7) |- (xor.west);
  \draw[->] (B) -| (2.6,1.3) |- (xor.west);
  \draw[->] (A) -| (2.4,0.7) |- (and.west);
  \draw[->] (B) -| (2.6,0.3) |- (and.west);
  \draw[->] (xor.east) -- ++(1,0) node[right] {$S$};
  \draw[->] (and.east) -- ++(1,0) node[right] {$C$};
\end{tikzpicture}
\end{center}

\medskip\hrule\vspace{4pt}

% -- 1.4 Tableau de Karnaugh -------------------------------------------
\textbf{1.4 Simplification par tableau de Karnaugh}

\textbf{Règles :} grouper les « 1 » adjacents par puissances de 2 (groupes de 2 → $-1$ var, de 4 → $-2$ var, de 8 → $-3$ var). Le tableau est \textbf{torique} : les bords sont adjacents.

\medskip\hrule\vspace{4pt}

% -- 1.5 Circuits combinatoires courants ------------------------------
\textbf{1.5 Circuits combinatoires courants}

\begin{center}
\renewcommand{\arraystretch}{1.2}
\begin{tabular}{lp{9cm}}
\toprule
\textbf{Circuit} & \textbf{Fonction} \\
\midrule
Multiplexeur $2^n \to 1$ & Sélectionne 1 entrée parmi $2^n$ selon $n$ bits de sélection \\
Décodeur $n \to 2^n$ & Active une sortie parmi $2^n$ (adressage mémoire) \\
Additionneur complet & $S = A \oplus B \oplus C_{in}$ ; $C_{out} = A\cdot B + C_{in}(A\oplus B)$ \\
Comparateur & Génère $A > B$, $A = B$, $A < B$ \\
\bottomrule
\end{tabular}
\end{center}

\medskip\hrule\vspace{4pt}

% -- 1.6 Codage numerique ----------------------------------------------
\textbf{1.6 Codages numériques}

\begin{itemize}
  \item \textbf{Binaire naturel :} $N_{10} = b_{n-1}\cdot 2^{n-1} + \cdots + b_1 \cdot 2 + b_0$
  \item \textbf{BCD :} chaque décimale codée sur 4 bits (ex. : $9 = 1001_2$)
  \item \textbf{Code Gray :} un seul bit change entre deux mots consécutifs - codeurs angulaires
  \item \textbf{Hexadécimal :} base 16 ; chiffres 0-9 puis A→F ; $0\text{x}45 = 69_{10}$
\end{itemize}

\textit{Connexion sujets CAPET (2023) :} codeur incrémental 12 bits, résolution $0{,}15$ mm/impulsion ; 8 bits de poids fort retenus (résolution effective $0{,}15 \times 2^4 = 2{,}4$ mm).

\end{tcolorbox}

% -----------------------------------------------------------------------
%  2. FORMULAIRE SYNTHETIQUE
% -----------------------------------------------------------------------

\begin{tcolorbox}[
  ficheSection=colAuto,
  title={2. Formulaire synthétique},
]
\small
\begin{multicols}{2}
\renewcommand{\arraystretch}{1.45}
\begin{tabular}{@{}p{5cm}l@{}}
\toprule
\textbf{Identité / formule} & \textbf{Nom} \\
\midrule
$\overline{A \cdot B} = \overline{A} + \overline{B}$ & De Morgan AND \\
$\overline{A + B} = \overline{A} \cdot \overline{B}$ & De Morgan OR \\
$A(A+B) = A$ & Absorption \\
$A+A\cdot B = A$ & Absorption \\
$A \oplus B = A\overline{B} + \overline{A}B$ & XOR développé \\
\bottomrule
\end{tabular}

\columnbreak

\renewcommand{\arraystretch}{1.45}
\begin{tabular}{@{}p{5cm}l@{}}
\toprule
\textbf{Formule} & \textbf{Usage} \\
\midrule
$S = A \oplus B \oplus C_{in}$ & Additionneur (somme) \\
$C_{out} = AB + C_{in}(A\oplus B)$ & Additionneur (retenue) \\
$n$ variables → $2^n$ lignes & Table vérité \\
$q = U_{PE}/2^n$ & Quantum CAN \\
$f_e \geq 2\,f_{max}$ & Shannon \\
\bottomrule
\end{tabular}
\end{multicols}
\normalsize
\end{tcolorbox}

% -----------------------------------------------------------------------
%  3. QUESTIONS TYPES
% -----------------------------------------------------------------------

\begin{tcolorbox}[
  ficheSection=colAuto,
  title={3. Questions types - extraites des sujets réels},
]

% -- QT1 -------------------------------------------------------------

\begin{tcolorbox}[questiontype,
  title={QT1 - Codage binaire / hexadécimal \hfill\normalfont\footnotesize E2 disciplinaire 2023 (D2-D3)}
]
\textbf{Énoncé.}
Un codeur incrémental génère une impulsion tous les $0{,}15$ mm. La distance maximale à coder est $680$ mm. Quel est le nombre minimal de bits nécessaires ? Le registre ne retient que les 8 bits de poids fort. Quelle est la résolution effective de position ?

\tcbline

\textbf{Corrigé.}

\textbf{Étape 1 — Nombre de bits.} Nombre de positions : $N = 680 / 0{,}15 = 4\,533$. Nombre de bits : $2^{12} = 4\,096 < 4\,533 < 8\,192 = 2^{13}$ - \textbf{13 bits}.

\textbf{Étape 2 — Résolution effective.} Les 8 bits de poids fort correspondent à ignorer les 5 bits de poids faible :
résolution effective $= 0{,}15 \times 2^5 = 0{,}15 \times 32 = \SI{4{,}8}{\milli\metre}$.

\textbf{Conclusion.} La troncature à 8 bits dégrade la résolution de \SI{0{,}15}{\milli\metre} à \SI{4{,}8}{\milli\metre} : ce compromis mémoire/précision doit être validé vis-à-vis de la tolérance mécanique exigée par le cahier des charges avant d'être retenu.
\end{tcolorbox}

\vspace{4pt}

% -- QT2 -------------------------------------------------------------

\begin{tcolorbox}[questiontype,
  title={QT2 - Application de De Morgan \hfill\normalfont\footnotesize Cours programme}
]
\textbf{Énoncé.}
Simplifier l'expression $S = \overline{\overline{A} \cdot \overline{B}}$ et identifier la porte correspondante. Puis simplifier $F = \overline{A + B + C}$.

\tcbline

\textbf{Corrigé.}

\textbf{Étape 1.} $S = \overline{\overline{A} \cdot \overline{B}} = \overline{\overline{A}} + \overline{\overline{B}} = A + B$ - \textbf{porte OR} réalisée par un NOR suivi d'un NON.

\textbf{Étape 2.} $F = \overline{A + B + C} = \overline{A} \cdot \overline{B} \cdot \overline{C}$ (De Morgan généralisé) - NOR à 3 entrées.

\textbf{Conclusion.} Ces deux simplifications confirment qu'une fonction OR (ou NOR à $n$ entrées) peut toujours se reconstruire à partir de portes NAND/NOR seules, ce qui justifie leur qualification de portes universelles.
\end{tcolorbox}

\vspace{4pt}

% -- QT3 -------------------------------------------------------------

\begin{tcolorbox}[questiontype,
  title={QT3 - Conversion hexadécimal / décimal \hfill\normalfont\footnotesize E2 disciplinaire 2022 (Q.17)}
]
\textbf{Énoncé.}
Un registre Siemens S7-1500 contient la valeur $\text{0x}2A$ (hexadécimal). Donner la valeur décimale. Un masquage par ET avec $\text{0x}0F$ est appliqué. Que conserve-t-on ?

\tcbline

\textbf{Corrigé.}

\textbf{Étape 1 — Conversion.} $\text{0x}2A = 2 \times 16 + 10 = 42_{10}$ ; en binaire : $0010\,1010_2$.

\textbf{Étape 2 — Masquage.} Masque $\text{0x}0F = 0000\,1111_2$. Le ET logique donne $0000\,1010_2 = 10_{10}$.

\textbf{Conclusion.} Le masque $\text{0x}0F$ isole les 4 bits de poids faible (le quartet de droite) et met à zéro les 4 bits de poids fort : c'est une opération classique d'extraction d'un champ dans un registre, très utilisée en programmation d'automate.
\end{tcolorbox}

\end{tcolorbox}

% -----------------------------------------------------------------------
%  4. ATTENTES DU JURY
% -----------------------------------------------------------------------

\begin{tcolorbox}[
  jury,
  title={4. Attentes du jury - Erreurs récurrentes et points de vigilance},
]

\begin{enumerate}[label=\textbf{\arabic*.}, leftmargin=*, itemsep=5pt]

  \item \textbf{De Morgan - le piège classique.} $\overline{A + B} = \overline{A} \cdot \overline{B}$ (ET), \textbf{pas} $\overline{A} + \overline{B}$. Toujours vérifier sur un cas numérique ($A=1$, $B=0$).

  \item \textbf{Tableau de Karnaugh torique.} Les cases des bords gauche et droit sont adjacentes, de même pour haut et bas. Oublier de boucler est une perte de simplification.

  \item \textbf{Confusion bit / octet.} Le jury 2023 signale : "beaucoup d'amalgames entre bit et octet". 1 octet = 8 bits ; $2^8 = 256$ valeurs.

  \item \textbf{XOR vs OR inclusif.} $A \oplus B = 1$ ssi $A \neq B$ (OR \emph{exclusif}). Le OR standard est $A = 1$ \textbf{ou} $B = 1$ \textbf{ou} les deux.

  \item \textbf{Codage binaire signé.} Pour les registres d'automate sur 16 bits signés (Siemens S7), le MSB est le bit de signe (complément à 2). Plage : $-32\,768$ à $+32\,767$.

  \item \textbf{Contexte CAPET.} La logique combinatoire pure (Karnaugh, SOP/POS) est peu présente dans les épreuves écrites. Maîtriser \emph{surtout} le codage binaire, hexadécimal, et les opérations de masquage/décalage qui apparaissent régulièrement.

\end{enumerate}

\medskip
\begin{tcolorbox}[
  colback=colJuryFrame!12!colJury,
  colframe=colJuryFrame,
  left=4pt, right=4pt, top=2pt, bottom=2pt,
  fontupper=\small\itshape
]
Rapport jury 2023 : «~La cinquième partie permettait d'aborder l'électronique numérique et la communication par liaison série. Le jury constate beaucoup d'amalgames entre bit et octet et invite les candidats à revoir le principe d'une liaison série.~»
\end{tcolorbox}

\end{tcolorbox}

% -----------------------------------------------------------------------
%  PIED - PRESENCE DANS LES SUJETS
% -----------------------------------------------------------------------

\begin{tcolorbox}[
  ficheSection=colPresence,
  title={Présence dans les sujets analysés (2013-2026) - logique numérique et codage},
]
\begin{center}
\footnotesize\renewcommand{\arraystretch}{1.15}
\begin{tabular}{lccccccccccccc}
\toprule
\textbf{Année} & 13 & 14 & 15 & 16 & 17 & 18 & 19 & 20 & 21 & 22 & 23 & 24 & 25 \\
\midrule
Codage binaire/hex & - & - & - & - & - & - & - & - & - & $\checkmark$ & $\checkmark$ & $\checkmark$ & - \\
Masquage / registre & - & - & - & - & - & - & - & - & - & $\checkmark$ & - & - & - \\
Logique comb. pure & - & - & - & - & - & - & - & - & - & - & - & - & - \\
\bottomrule
\end{tabular}
\end{center}
\smallskip
\footnotesize\textbf{Sous-thèmes fréquents :}
codage binaire - résolution/précision - conversion hexadécimal - masquage ET/OU - nombre de bits pour coder une position - conversion décimal/binaire/hexa

\end{tcolorbox}

\CAPETENDFILE
